\input{../tex-inputs/latex-leseansicht-vorspann}

\section[Gustav Schwarzkopf an Arthur Schnitzler, 16. 7. 1916]{L04252 Gustav Schwarzkopf an Arthur Schnitzler, 16. 7. 1916}
\nopagebreak\mylabel{L04252v}
\rehead{ }\normalsize\beginnumbering\briefempfaengerindex{Schnitzler, Arthur@\textsc{Schnitzler, Arthur}!zzzSchwarzkopf, Gustav@\emph{von Gustav Schwarzkopf}!1916-07-161@{16. 7. 1916}|(be}
\toendnotes[C]{\smallbreak\pagebreak[2]}
\correspDesc{Versand  durch Gustav Schwarzkopf am 16. 7. 1916 in Marienbad
\newline{}Übermittlung  am 17. 7. 1916 in Marienbad
\newline{}Erhalt  durch Arthur Schnitzler im Zeitraum [17. 7. 1916 – 21. 7. 1916?] in Altaussee}\toendnotes[C]{\smallbreak}
\Standort{CUL, Schnitzler, B 96a.}
\physDesc{Bildpostkarte, 529 Zeichen
\newline{}Handschrift: schwarze Tinte, deutsche Kurrent
\newline{}Versand: Stempel: »\nobreak{}\oindex{Marienbad@\textbf{Marienbad}|pwk}Marien\textcolor{gray}{b}{[}ad{]}, 17. VII. 16, VI\nobreak{}«.  }\toendnotes[C]{\smallbreak}\pstart{}{\pb}Herrn\pend{}\pstart{}\textsc{D\textsuperscript{r} Arthur
                  Schnitzler}\pend{}\pstart{}\textsc{Alt-Aussee\oindex{Altaussee@\textbf{Altaussee}, \emph{Verwaltungsgebiet}|pw}}\pend{}{\bigskip}
\pstart
           \noindent{}\centering{}{\pb}\textcolor{gray}{\textbf{Marienbad. Rübezahl\oindex{Hotel Rübezahl@\textbf{Hotel Rübezahl}, \emph{Hotel}|pw}.}}\pend
           \vspace{1em}
\pstart
           Villa Helgoland\oindex{Villa Helgoland@\textbf{Villa Helgoland}, \emph{Beherbergungsgebäude}|pw}. Z\textsuperscript{r} 18.\pend
           
\pstart
           (Waldquellzeile\oindex{Třebízského ulice [Marienbad]@\textbf{Třebízského ulice [Marienbad]}, \emph{Meerenge}|pw})\pend
           
\pstart
           \raggedleft{}{\pb}16. VII 16\pend
           
\pstart{}Lieber Arthur!\pend\vspace{0.5em}
\pstart
           Seit Freitag Abend hier. Zu meiner Überraſchung. Es iſt hier{ }ſehr{ }ſchön,
               es gibt ein Glück, das mit Se{\geminationm}eln eine entfernte
               Ähnlichkeit hat und es regnet hier genau{ }ſo tüchtig und ausdauernd wie in \textsc{Ischl\oindex{Bad Ischl@\textbf{Bad Ischl}|pw}} oder \textsc{Aussee\oindex{Bad Aussee@\textbf{Bad Aussee}, \emph{Hauptstadt}|pw}}.\pend
           
\pstart
           Ich will noch etwa 14 Tage hier bleiben –{ }ſo lange eben »die vorhandenen Mittel
               reichen«. – \uline{Wer} hat bei dem Konzert in Ihrer Villa\oindex{Villa Annerl@\textbf{Villa Annerl}, \emph{Gebäude}|pwv} zugehört? Die \label{K_L04252-1v}\edtext{\uline{Weltweiſe}}{\lemma{\textnormal{\emph{Weltweise}}}\Cendnote{\textnormal{{XXXX ref} XXXX 7.7.1916}}}\label{K_L04252-1}? Ich ko{\geminationn}te es
               nicht entziffern.\pend
           
\pstart
           Ihre Frau\pwindex{Schnitzler, Olga 17.\,1.\,1882 Wien – 13.\,1.\,1970 Lugano@\textsc{Schnitzler, Olga} (17.\,1.\,1882 Wien – 13.\,1.\,1970 Lugano), \emph{Schauspielerin, Sängerin}|pwv} und{\\[\baselineskip]}Sie
               herzlichſt grüßend Ihr \spacefill\mbox{Gustav}\pend
           \leftskip=0em{}\selectlanguage{ngerman}\endnumbering\briefempfaengerindex{Schnitzler, Arthur@\textsc{Schnitzler, Arthur}!zzzSchwarzkopf, Gustav@\emph{von Gustav Schwarzkopf}!1916-07-161@{16. 7. 1916}|)be}\mylabel{L04252h}
\begin{anhang}
\end{anhang}\newcommand{\dateiname}{L04252}\newcommand{\titel}{Gustav Schwarzkopf an Arthur Schnitzler, 16. 7. 1916}\newcommand{\editorInnen}{Herausgegeben von Jahnke, SelmaMüller, Martin Anton}\input{../tex-inputs/latex-leseansicht-abspann}
